\section{本讲小结}
动力学题要联立电路方程、受力方程和初态；能量用于检查或简化，不能替代缺失的运动条件。纯电阻、接电容与几何变长分别产生不同反馈，不能共用一个磁阻力公式。

多杆系统先识别相对速度和可用的不变量；碰撞与电磁耗散分别计算。跨磁区、跨轨宽台阶和击穿都要分段，指数趋近、准静态末段与局部涡流估算则要说明其精度和适用范围。
